\section{Bối cảnh pháp lý và đặc thù DNNVV}

\subsection{Doanh nghiệp nhỏ và vừa trong hệ thống pháp luật Việt Nam}

Luật Hỗ trợ doanh nghiệp nhỏ và vừa số 04/2017/QH14 xác định DNNVV theo tiêu
chí quy mô (lao động, doanh thu, nguồn vốn) gắn với ngành nghề, và gắn một số
chính sách hỗ trợ~\cite{quochoi2017}. Đồ án không tái lập toàn bộ tiêu chí phân
loại để cấp ``giấy chứng nhận DNNVV''. Hồ sơ tổ chức trong hệ thống chỉ lưu các
tham số đủ để sinh lịch và mặc định chủ thể khi hỏi đáp: tên, mã số thuế, số
lao động, kỳ kê khai GTGT.

Điểm thực tiễn: nghĩa vụ thuế, bảo hiểm, báo cáo lao động và hóa đơn không biến
mất khi quy mô nhỏ. Bộ phận pháp chế riêng thì thường không có. Người dùng thật
của đồ án vì thế là chủ doanh nghiệp, kế toán dịch vụ hoặc nhân sự kiêm nhiệm
--- đối tượng hỏi bằng lời nói hàng ngày hơn là bằng số hiệu thông tư.

\subsection{Hệ thống văn bản quy phạm và thứ bậc áp dụng}

Văn bản do Quốc hội ban hành (luật, bộ luật) đặt nguyên tắc. Nghị định của
Chính phủ quy định chi tiết. Thông tư của bộ hướng dẫn thực hiện. Câu hỏi thao
tác (``nộp tờ khai ngày nào'', ``hóa đơn gửi khi nào'') thường rơi xuống nghị
định và thông tư, trong khi câu hỏi nguyên tắc (``có được đơn phương chấm dứt
không'') rơi ở luật. Corpus đồ án cố ý chứa cả ba cấp trong các lĩnh vực đã chọn,
tránh tình trạng trả lời đúng luật khung nhưng sai bước kê khai.

Hiệu lực theo thời gian là bài toán pháp điển riêng: văn bản mới thay văn bản
cũ, điều khoản chuyển tiếp, văn bản hết hiệu lực còn nằm trên cổng. Đồ án lưu
theo số hiệu, chưa gắn lịch sử hiệu lực từng Điều. Khi kho có hai luật cùng tên
khác số (ví dụ quản lý thuế, thuế GTGT), cả hai đều tìm được; người dùng và
người quản trị phải tự chọn bản đang áp dụng. Hạn chế này được nêu lại ở
Chương~5, không giấu trong phần ``đã hoàn thiện''.

\subsection{Rủi ro tuân thủ đối với người không chuyên}

Chậm nộp, sai tờ khai, thiếu báo cáo lao động có thể kéo theo tiền chậm nộp và
xử phạt hành chính. Chatbot không nguồn tạo rủi ro kiểu khác: câu tiếng Việt
trôi, kèm số Điều nghe quen, khiến người dùng tin hơn kết quả tìm từ khóa thô.
Đồ án đặt yêu cầu tối thiểu: mọi căn cứ đưa ra UI phải mở được về một Điều trong
kho. Yêu cầu này không biến hệ thống thành tư vấn pháp lý có trách nhiệm nghề
nghiệp. Người dùng vẫn phải đối chiếu văn bản gốc và, khi tranh chấp, nhờ tổ
chức hành nghề.

\subsection{Ranh giới trợ lý AI và dịch vụ pháp lý}

Hệ thống hỗ trợ tra cứu, nhắc việc và chỉ ra chỗ hợp đồng có thể lệch so với
Điều đã truy hồi. Hệ thống không đại diện tố tụng, không soạn hồ sơ nộp cơ quan
thuế, không cam kết ``thắng kiện''. Câu trả lời mang tính hỗ trợ thông tin trên
corpus đóng. Ranh giới này vừa là yêu cầu đạo đức nghề, vừa là phạm vi đồ án
ngành: thời gian không đủ để xây phân hệ chữ ký số hay kênh nộp thật.

\section{Bài toán tin học đặt ra từ hiện trạng}

Từ hiện trạng trên, bài toán không phải ``cài chatbot''. Ba lớp việc phải làm
cùng lúc:

\begin{enumerate}
    \item \textbf{Lớp dữ liệu.} Chuẩn hóa \soLuat{} văn bản thành Điều có siêu
    dữ liệu (số hiệu, chương, viện dẫn), nhập PostgreSQL, lập chỉ mục không dấu,
    nhúng vector.
    \item \textbf{Lớp suy diễn hạn chế.} Pipeline RAG lai, hậu kiểm citation,
    lịch theo rule --- suy diễn nằm ở chỗ ghép nguồn, không ở chỗ mô hình tự
    đặt ra hạn nộp.
    \item \textbf{Lớp phần mềm.} Đa người dùng, JWT, cách ly tổ chức, SSE, Docker.
    Không có lớp này thì pipeline chỉ chạy được trên máy phát triển.
\end{enumerate}

Báo cáo lần lượt trả lời ba lớp: Chương~2 cho kỹ thuật truy hồi và công nghệ;
Chương~3 cho phân tích thiết kế; Chương~4 cho bằng chứng chạy được.
